\documentclass[11pt]{article}

\usepackage[margin=1in]{geometry}
\usepackage{amsmath,amssymb,amsthm,mathtools}
\usepackage{enumitem}
\usepackage{hyperref}
\usepackage{mathrsfs}
\usepackage{microtype}

\hypersetup{
  colorlinks=true,
  linkcolor=blue,
  citecolor=blue,
  urlcolor=blue
}

\newtheorem{theorem}{Theorem}[section]
\newtheorem{proposition}[theorem]{Proposition}
\newtheorem{lemma}[theorem]{Lemma}
\newtheorem{corollary}[theorem]{Corollary}
\newtheorem{condition}[theorem]{Condition}
\theoremstyle{definition}
\newtheorem{definition}[theorem]{Definition}
\newtheorem{remark}[theorem]{Remark}

\newcommand{\E}{\mathbb{E}}
\newcommand{\Prob}{\mathbb{P}}
\newcommand{\R}{\mathbb{R}}
\newcommand{\1}{\mathbf{1}}
\DeclareMathOperator{\spn}{sp}
\DeclareMathOperator{\supp}{supp}

\title{Uniform $\varepsilon$-Equilibria in Finite Stochastic Games\\under a Global Irreducibility Condition}
\author{Anonymous draft}
\date{v27}

\begin{document}
\maketitle

\begin{abstract}
We prove a short uniform-equilibrium theorem for finite stochastic games under a strong game-level communication hypothesis. The hypothesis is the following global irreducibility condition: for every player $i$, every pure stationary strategy of player $i$, and every mixed stationary profile of the other players, the induced Markov chain on the original state space $S$ is irreducible. Under this condition, each unilateral decision problem is a finite unichain average-reward Markov decision process on $S$. We then study the normal-form game in which players choose stationary mixed strategies and receive their ergodic average payoffs. The payoff functions of this meta-game are continuous; the best-response correspondence has nonempty compact convex values; and the Kakutani-Fan-Glicksberg fixed-point theorem yields a stationary average-reward Nash equilibrium $x^*$. For every player, the bias function of the induced average-reward MDP against $x^*_{-i}$ gives a supermartingale that bounds the payoff of \emph{every} unilateral deviation, including arbitrary history-dependent behavioral deviations. Because $x_i^*$ is stationary average-reward optimal, the compliant payoff is within $\spn(h_i^*)/T$ of the optimal gain at horizon $T$. Hence the regret is at most $2\max_i \spn(h_i^*)/T$, so $x^*$ is a uniform $\varepsilon$-equilibrium for all sufficiently long finite horizons. The proof works directly on the original game and uses no tree decomposition, local routing, or auxiliary state-space expansion.
\end{abstract}

\section{Introduction}

Finite stochastic games, introduced by Shapley \cite{Shapley1953}, form the basic finite-state model of dynamic strategic interaction with publicly observed states. In the zero-sum two-player case, the long-run average theory is classical: Mertens and Neyman \cite{MertensNeyman1981} proved the existence of the uniform value. In the nonzero-sum two-player case, Vieille \cite{Vieille2000a,Vieille2000b} proved the existence of uniform equilibrium payoffs. For general multiplayer finite stochastic games, however, the existence of uniform equilibria remains open; see the surveys of Solan and Vieille \cite{SolanVieille2015}, Solan \cite{Solan2022}, and the current open-problem list maintained by the Game Theory Society \cite{OpenProblems2026}.

This note records a simple positive result under a strong game-level communication hypothesis. The condition is imposed directly on the original state space $S$ and directly on unilateral stationary deviations. We do \emph{not} expand the state space by any auxiliary tree, routing device, or controller memory. Under the hypothesis, the entire state space is one communicating class for every unilateral pure stationary deviation, and therefore each player faces a unichain average-reward MDP on $S$ when the opponents hold a stationary mixed profile fixed. This is exactly the setting in which the standard gain-bias machinery applies cleanly.

The condition that we use is the following one.

\begin{quote}
For every player $i$, every pure stationary strategy $f_i$ of player $i$, and every mixed stationary profile $x_{-i}$ of the other players, the Markov chain on $S$ induced by $(f_i,x_{-i})$ is irreducible.
\end{quote}

This is a genuinely game-level condition. Once it holds, there is no need to decompose the game into smaller communicating pieces. Indeed, any auxiliary expansion of the state space can create recurrent classes that are absent in the original game, whereas the argument below derives all estimates directly on $S$.

Our main result is the following.

\begin{theorem}[Main theorem]\label{thm:intro-main}
Let $G$ be a finite stochastic game satisfying Condition~\ref{cond:A}. Then for every $\varepsilon>0$ there exist a stationary mixed strategy profile $x^*$ and an integer $T_0$ such that for every initial state $s\in S$, every horizon $T\ge T_0$, every player $i$, and every behavioral deviation $\tau_i$ of player $i$,
\[
\gamma_i^T(s,x^*)\;\ge\;\gamma_i^T\bigl(s,(\tau_i,x^*_{-i})\bigr)-\varepsilon.
\]
Hence $x^*$ is a uniform $\varepsilon$-equilibrium.

More precisely, if $h_i^*$ is an optimal bias for player $i$ against $x^*_{-i}$, then one may take
\[
T_0\;=\;\left\lceil \frac{2\max_{i}\spn(h_i^*)}{\varepsilon}\right\rceil.
\]
\end{theorem}

The proof is a five-step chain.
\begin{enumerate}[label=\textup{(\arabic*)},leftmargin=2.8em]
\item Define the average-reward meta-game in which players choose stationary mixed strategies and receive their ergodic average rewards.
\item Show that every stationary mixed profile induces an irreducible Markov chain, hence an invariant distribution and an average payoff vector $g(x)$; then prove that $g$ is continuous.
\item Fix opponents' stationary profile. Condition~\ref{cond:A} turns the unilateral optimization problem of player $i$ into a finite unichain average-reward MDP on $S$, with optimal gain $g_i^*$ and bias $h_i^*$.
\item Use the Bellman inequality to build a supermartingale that bounds the payoff of \emph{every} unilateral deviation, including arbitrary history-dependent randomized deviations.
\item Use the fixed-point theorem of Kakutani-Fan-Glicksberg to obtain a stationary average-reward Nash equilibrium $x^*$ of the meta-game, and combine the deviation upper bound with the finite-horizon identity for optimal stationary policies.
\end{enumerate}

The conclusion is quantitative: the regret at horizon $T$ is $O(1/T)$. In particular, the stationary profile does not depend on the horizon. The game-level irreducibility hypothesis is doing all the structural work. It suppresses multichain behavior at the level of the original game and allows the entire proof to be written with standard average-reward MDP theory and a single fixed-point argument.

The paper is organized as follows. Section~\ref{sec:model} introduces the model and proves that Condition~\ref{cond:A} implies irreducibility for every stationary mixed profile. Section~\ref{sec:mdp} records the average-reward MDP facts that we use, including the supermartingale bound for arbitrary deviations. Section~\ref{sec:metagame} defines the average-reward meta-game, proves continuity of the payoff map, and obtains a stationary average-reward equilibrium by a Kakutani-Fan-Glicksberg argument. Section~\ref{sec:proof-main} proves Theorem~\ref{thm:intro-main}. Section~\ref{sec:discussion} closes with a short discussion.

\section{Finite stochastic games and the global irreducibility condition}\label{sec:model}

\subsection{Game model}

A finite stochastic game is given by
\[
G=(N,S,(A_i)_{i\in N},u,P),
\]
where:
\begin{itemize}[leftmargin=2.3em]
\item $N=\{1,\dots,n\}$ is the finite set of players,
\item $S$ is a finite nonempty set of states,
\item for each player $i$ and each state $s\in S$, the action set $A_i(s)$ is finite and nonempty,
\item for each player $i$, the stage payoff function satisfies
\[
 u_i(s,a)\in[0,1]\qquad\text{for all }s\in S,\ a=(a_j)_{j\in N}\in A(s):=\prod_{j\in N}A_j(s),
\]
\item $P(\cdot\mid s,a)\in\Delta(S)$ is the transition law from state $s$ under action profile $a$.
\end{itemize}

At stage $t$, the current state $s_t$ is publicly observed, the players simultaneously choose an action profile $a_t=(a_{1,t},\dots,a_{n,t})\in A(s_t)$, players receive payoffs $u_i(s_t,a_t)$, and the next state is drawn according to $P(\cdot\mid s_t,a_t)$.

A \emph{behavioral strategy} of player $i$ is a rule $\tau_i$ assigning to every finite public history
\[
h_t=(s_1,a_1,s_2,a_2,\dots,s_t)
\]
a probability distribution $\tau_i(\cdot\mid h_t)\in\Delta(A_i(s_t))$. A strategy profile is denoted by $\sigma=(\sigma_1,\dots,\sigma_n)$.

For $T\ge1$, the expected $T$-stage average payoff of player $i$ from initial state $s$ under profile $\sigma$ is
\[
\gamma_i^T(s,\sigma)
=\E_s^{\sigma}\!\left[\frac1T\sum_{t=1}^{T}u_i(s_t,a_t)\right].
\]

\begin{definition}
A strategy profile $\sigma$ is a \emph{uniform $\varepsilon$-equilibrium} if there exists $T_0\in\mathbb N$ such that for every horizon $T\ge T_0$, every initial state $s\in S$, every player $i$, and every behavioral deviation $\tau_i$ of player $i$,
\[
\gamma_i^T(s,\sigma)\ge\gamma_i^T\bigl(s,(\tau_i,\sigma_{-i})\bigr)-\varepsilon.
\]
\end{definition}

We will use the following stationary strategy spaces.
\[
X_i:=\prod_{s\in S}\Delta(A_i(s)),\qquad X:=\prod_{i\in N}X_i.
\]
An element $x_i\in X_i$ is a \emph{stationary mixed strategy}. We write $x_i(a_i\mid s)$ for the probability assigned to action $a_i\in A_i(s)$. A \emph{pure stationary strategy} of player $i$ is a map $f_i:S\to\bigcup_{s}A_i(s)$ such that $f_i(s)\in A_i(s)$ for all $s$; the finite set of such strategies is denoted by
\[
F_i:=\prod_{s\in S}A_i(s).
\]

Given $x=(x_1,\dots,x_n)\in X$, define its stagewise product distribution by
\[
x(a\mid s):=\prod_{j\in N}x_j(a_j\mid s),\qquad a\in A(s),\ s\in S.
\]
Then the induced transition matrix on $S$ is
\[
P^x(s,s'):=\sum_{a\in A(s)}x(a\mid s)P(s'\mid s,a),\qquad s,s'\in S.
\]
For a pure stationary strategy $f_i\in F_i$ and an opponents' stationary mixed profile $x_{-i}\in X_{-i}$, we similarly write
\[
P^{f_i,x_{-i}}(s,s'):=\sum_{a_{-i}\in A_{-i}(s)}x_{-i}(a_{-i}\mid s)
P\bigl(s'\mid s,f_i(s),a_{-i}\bigr).
\]

\subsection{Global irreducibility}

\begin{condition}[Global irreducibility]\label{cond:A}
For every player $i\in N$, every pure stationary strategy $f_i\in F_i$, and every stationary mixed profile $x_{-i}\in X_{-i}$ of the other players, the Markov chain on $S$ with transition matrix $P^{f_i,x_{-i}}$ is irreducible.
\end{condition}

\begin{remark}
Condition~\ref{cond:A} is imposed on the original game and on the original state space $S$. A simple sufficient condition is a Doeblin-type lower bound: if there exist $\alpha>0$ and a probability measure $q$ with full support on $S$ such that
\[
P(\cdot\mid s,a)\ge \alpha q(\cdot)\qquad\text{for every }s\in S\text{ and }a\in A(s),
\]
then Condition~\ref{cond:A} holds automatically. The point of Condition~\ref{cond:A}, however, is that it does \emph{not} require such a uniform minorization. It only requires irreducibility of the full chain on $S$ under every unilateral pure stationary deviation against stationary mixed opponents.
\end{remark}

The first basic consequence is that mixed stationary strategies do not destroy irreducibility.

\begin{proposition}\label{prop:mixed-irreducible}
Assume Condition~\ref{cond:A}. Then for every stationary mixed profile $x\in X$, the Markov chain on $S$ with transition matrix $P^x$ is irreducible.
\end{proposition}

\begin{proof}
Fix $x\in X$ and fix one player $i\in N$. For each pure stationary strategy $f_i\in F_i$, define
\[
\lambda_{f_i}(x_i):=\prod_{s\in S}x_i\bigl(f_i(s)\mid s\bigr).
\]
Because the action choices at different states factor, we have
\[
\sum_{f_i\in F_i}\lambda_{f_i}(x_i)=1.
\]
Moreover, a direct expansion over pure stationary strategies gives the matrix identity
\begin{equation}\label{eq:convex-combination}
P^x=\sum_{f_i\in F_i}\lambda_{f_i}(x_i)\,P^{f_i,x_{-i}}.
\end{equation}
Indeed, for each pair of states $s,s'$, the $(s,s')$-entry of the right-hand side equals
\begin{align*}
\sum_{f_i\in F_i}\lambda_{f_i}(x_i)P^{f_i,x_{-i}}(s,s')
&=\sum_{f_i\in F_i}\left(\prod_{z\in S}x_i(f_i(z)\mid z)\right)
\sum_{a_{-i}}x_{-i}(a_{-i}\mid s)P(s'\mid s,f_i(s),a_{-i})\\
&=\sum_{a_i\in A_i(s)}x_i(a_i\mid s)\sum_{a_{-i}}x_{-i}(a_{-i}\mid s)P(s'\mid s,a_i,a_{-i})\\
&=P^x(s,s').
\end{align*}
Choose $f_i^\sharp\in F_i$ such that $x_i(f_i^\sharp(s)\mid s)>0$ for every state $s$. Then $\lambda_{f_i^\sharp}(x_i)>0$. By Condition~\ref{cond:A}, the chain $P^{f_i^\sharp,x_{-i}}$ is irreducible. Therefore, for every pair of states $p,q\in S$, there exists $m\ge1$ such that
\[
\bigl(P^{f_i^\sharp,x_{-i}}\bigr)^m(p,q)>0.
\]
From \eqref{eq:convex-combination} we have the entrywise inequality
\[
P^x\ge \lambda_{f_i^\sharp}(x_i)\,P^{f_i^\sharp,x_{-i}},
\]
whence, by monotonicity of matrix multiplication for nonnegative matrices,
\[
(P^x)^m\ge \lambda_{f_i^\sharp}(x_i)^m\bigl(P^{f_i^\sharp,x_{-i}}\bigr)^m.
\]
In particular,
\[
(P^x)^m(p,q)\ge \lambda_{f_i^\sharp}(x_i)^m\bigl(P^{f_i^\sharp,x_{-i}}\bigr)^m(p,q)>0.
\]
Thus $P^x$ is irreducible.
\end{proof}

For the unilateral optimization problems that will appear later, this immediately implies a unichain property.

\begin{corollary}\label{cor:unichain-from-A}
Assume Condition~\ref{cond:A}, fix a player $i$, and fix an opponents' stationary mixed profile $x_{-i}\in X_{-i}$. Consider the MDP faced by player $i$ on state space $S$, with action set $A_i(s)$ at state $s$, transition kernel
\[
Q_i^{x_{-i}}(s'\mid s,a_i):=\sum_{a_{-i}\in A_{-i}(s)}x_{-i}(a_{-i}\mid s)P(s'\mid s,a_i,a_{-i}),
\]
and reward
\[
r_i^{x_{-i}}(s,a_i):=\sum_{a_{-i}\in A_{-i}(s)}x_{-i}(a_{-i}\mid s)u_i(s,a_i,a_{-i}).
\]
Then every pure stationary policy of this MDP induces an irreducible Markov chain on $S$. In particular, the MDP is unichain; indeed it is irreducible under every pure stationary policy.
\end{corollary}

\begin{proof}
A pure stationary policy of the MDP is exactly a pure stationary strategy $f_i\in F_i$ of player $i$ in the stochastic game. Its induced transition matrix is $P^{f_i,x_{-i}}$, which is irreducible by Condition~\ref{cond:A}.
\end{proof}

\section{Average-reward MDP facts under fixed opponents}\label{sec:mdp}

Fix a player $i\in N$ and an opponents' stationary mixed profile $x_{-i}\in X_{-i}$. The controlled process of player $i$ is the finite MDP from Corollary~\ref{cor:unichain-from-A}.

For a stationary mixed policy $\mu_i\in X_i$, its induced transition matrix is
\[
P^{\mu_i,x_{-i}}(s,s'):=\sum_{a_i\in A_i(s)}\mu_i(a_i\mid s)
Q_i^{x_{-i}}(s'\mid s,a_i),
\]
and its expected one-step reward in state $s$ is
\[
\bar r_i^{\mu_i,x_{-i}}(s):=\sum_{a_i\in A_i(s)}\mu_i(a_i\mid s)r_i^{x_{-i}}(s,a_i).
\]
By Proposition~\ref{prop:mixed-irreducible}, the chain $P^{\mu_i,x_{-i}}$ is irreducible, so it has a unique invariant distribution $\pi^{\mu_i,x_{-i}}$ and an average reward
\[
g_i(\mu_i,x_{-i})
:=\sum_{s\in S}\pi^{\mu_i,x_{-i}}(s)\,\bar r_i^{\mu_i,x_{-i}}(s),
\]
independent of the initial state.

\subsection{Gain and bias}

Fix a reference state $s^\dagger\in S$ once and for all. The following is standard finite-state average-reward MDP theory; see Puterman \cite[Chapter~8]{Puterman1994} or Neyman \cite[Section~3]{Neyman2003MarkovChains}.

\begin{theorem}[Average-reward optimality equation]\label{thm:aroe}
Assume Condition~\ref{cond:A}. For each player $i$ and each opponents' stationary mixed profile $x_{-i}\in X_{-i}$, there exist a scalar $g_i^*(x_{-i})\in\R$ and a function $h_i^{x_{-i}}:S\to\R$, unique under the normalization $h_i^{x_{-i}}(s^\dagger)=0$, such that
\begin{equation}\label{eq:bellman-equality}
g_i^*(x_{-i})+h_i^{x_{-i}}(s)
=\max_{a_i\in A_i(s)}\Bigl[r_i^{x_{-i}}(s,a_i)+\sum_{s'\in S}Q_i^{x_{-i}}(s'\mid s,a_i)h_i^{x_{-i}}(s')\Bigr]
\end{equation}
for every $s\in S$.

The number $g_i^*(x_{-i})$ is the optimal average reward of player $i$ against $x_{-i}$ over all admissible policies, including history-dependent randomized policies. Moreover, every pure stationary selector of the maximizers in \eqref{eq:bellman-equality} is average-reward optimal.
\end{theorem}

\begin{proof}
By Corollary~\ref{cor:unichain-from-A}, the induced MDP is finite and unichain. The existence of a unique normalized solution $(g_i^*(x_{-i}),h_i^{x_{-i}})$ to the average-reward optimality equation, together with the optimality of stationary maximizing selectors, is a standard theorem for finite unichain average-reward MDPs; see the references cited above.
\end{proof}

Write
\[
A_i^*(s;x_{-i})
:=\arg\max_{a_i\in A_i(s)}\Bigl[r_i^{x_{-i}}(s,a_i)+\sum_{s'\in S}Q_i^{x_{-i}}(s'\mid s,a_i)h_i^{x_{-i}}(s')\Bigr].
\]
Each set $A_i^*(s;x_{-i})$ is nonempty by finiteness.

\subsection{The supermartingale bound for arbitrary deviations}

The next proposition is the key estimate in the entire argument. It bounds the finite-horizon payoff of \emph{every} unilateral deviation, not merely stationary deviations.

\begin{proposition}[Supermartingale upper bound]\label{prop:supermartingale}
Assume Condition~\ref{cond:A}. Fix a player $i$ and a stationary mixed opponents' profile $x_{-i}\in X_{-i}$. Let $g_i^*=g_i^*(x_{-i})$ and $h_i=h_i^{x_{-i}}$ be the gain and normalized bias from Theorem~\ref{thm:aroe}. Then for every initial state $s\in S$, every horizon $T\ge1$, and every behavioral strategy $\tau_i$ of player $i$,
\begin{equation}\label{eq:deviation-upper}
\gamma_i^T\bigl(s,(\tau_i,x_{-i})\bigr)
\le g_i^*+\frac{\spn(h_i)}{T}.
\end{equation}
Equivalently, if $u_t:=u_i(s_t,a_t)$ and
\[
M_t:=h_i(s_t)+\sum_{k=1}^{t-1}(g_i^*-u_k),\qquad t\ge1,
\]
then $(M_t)_{t\ge1}$ is a supermartingale under $\Prob_s^{(\tau_i,x_{-i})}$.
\end{proposition}

\begin{proof}
Fix $s$, $T$, and $\tau_i$. Let $(\mathcal F_t)_{t\ge1}$ be a filtration containing the public history and the private randomization of player $i$, so that $s_t$ and the conditional law of $a_{i,t}$ are $\mathcal F_t$-measurable. Denote by $\mu_t(\cdot)$ the conditional distribution of $a_{i,t}$ given $\mathcal F_t$. Since the other players use the stationary mixed profile $x_{-i}$, conditional on $\mathcal F_t$ and on the event $\{a_{i,t}=a_i\}$, the expected reward and next-state law are exactly $r_i^{x_{-i}}(s_t,a_i)$ and $Q_i^{x_{-i}}(\cdot\mid s_t,a_i)$.

By Theorem~\ref{thm:aroe}, for every state $s$ and every action $a_i\in A_i(s)$,
\[
r_i^{x_{-i}}(s,a_i)+\sum_{s'}Q_i^{x_{-i}}(s'\mid s,a_i)h_i(s')\le g_i^*+h_i(s).
\]
Applying this inequality at the random state $s_t$ and averaging over the conditional action distribution $\mu_t$ yields
\begin{align*}
\E\bigl[u_t+h_i(s_{t+1})\mid\mathcal F_t\bigr]
&=\sum_{a_i}\mu_t(a_i)\left(r_i^{x_{-i}}(s_t,a_i)+\sum_{s'}Q_i^{x_{-i}}(s'\mid s_t,a_i)h_i(s')\right)\\
&\le g_i^*+h_i(s_t).
\end{align*}
Rearranging gives
\[
\E[M_{t+1}\mid \mathcal F_t]\le M_t,
\]
so $(M_t)$ is a supermartingale.

Taking expectations and using $M_1=h_i(s)$, we obtain
\[
\E[h_i(s_{T+1})]+Tg_i^*-\E\Bigl[\sum_{t=1}^{T}u_t\Bigr]\le h_i(s).
\]
Hence
\[
\E\Bigl[\sum_{t=1}^{T}u_t\Bigr]\le Tg_i^*+h_i(s)-\E[h_i(s_{T+1})].
\]
Because the difference between any two values of $h_i$ is bounded by its span,
\[
h_i(s)-\E[h_i(s_{T+1})]\le \max_{z\in S}h_i(z)-\min_{z\in S}h_i(z)=\spn(h_i).
\]
Dividing by $T$ yields \eqref{eq:deviation-upper}.
\end{proof}

\subsection{Optimal stationary responses and the exact finite-horizon identity}

The next result identifies the stationary optimal responses to $x_{-i}$ and gives the matching lower bound along such responses.

\begin{proposition}[Optimal stationary responses]\label{prop:optimal-stationary}
Assume Condition~\ref{cond:A}. Fix a player $i$ and a stationary mixed opponents' profile $x_{-i}\in X_{-i}$. Let $g_i^*=g_i^*(x_{-i})$ and $h_i=h_i^{x_{-i}}$ be as in Theorem~\ref{thm:aroe}.

\begin{enumerate}[label=\textup{(\alph*)},leftmargin=2.4em]
\item If a stationary mixed strategy $\mu_i\in X_i$ satisfies
\[
\supp \mu_i(\cdot\mid s)\subseteq A_i^*(s;x_{-i})\qquad\text{for every }s\in S,
\]
then for every initial state $s\in S$ and every horizon $T\ge1$,
\begin{equation}\label{eq:exact-identity}
\gamma_i^T\bigl(s,(\mu_i,x_{-i})\bigr)
= g_i^*+\frac{h_i(s)-\E_s^{(\mu_i,x_{-i})}[h_i(s_{T+1})]}{T}.
\end{equation}
In particular,
\begin{equation}\label{eq:exact-lower-bound}
\left|\gamma_i^T\bigl(s,(\mu_i,x_{-i})\bigr)-g_i^*\right|\le \frac{\spn(h_i)}{T},
\end{equation}
so $\mu_i$ is average-reward optimal.

\item Conversely, if a stationary mixed strategy $\mu_i\in X_i$ is average-reward optimal against $x_{-i}$, then
\[
\supp \mu_i(\cdot\mid s)\subseteq A_i^*(s;x_{-i})\qquad\text{for every }s\in S.
\]
Hence the stationary optimal responses are exactly the stationary mixed strategies whose support is contained statewise in the maximizing action sets.
\end{enumerate}
\end{proposition}

\begin{proof}
\textit{Part (a).} Let $\mu_i$ satisfy the support condition. Averaging the Bellman equality \eqref{eq:bellman-equality} over $\mu_i(\cdot\mid s)$ gives, for each state $s$,
\[
g_i^*+h_i(s)
=\bar r_i^{\mu_i,x_{-i}}(s)+\sum_{s'\in S}P^{\mu_i,x_{-i}}(s,s')h_i(s').
\]
Therefore, under the profile $(\mu_i,x_{-i})$,
\[
\E\bigl[u_t+h_i(s_{t+1})\mid s_t\bigr]=g_i^*+h_i(s_t).
\]
Taking full expectations and summing from $t=1$ to $T$ yields
\[
\E\Bigl[\sum_{t=1}^{T}u_t\Bigr]+\E[h_i(s_{T+1})]=Tg_i^*+h_i(s),
\]
which is exactly \eqref{eq:exact-identity}. The bound \eqref{eq:exact-lower-bound} follows immediately because the numerator is bounded in absolute value by $\spn(h_i)$. Thus $\mu_i$ attains the optimal average reward $g_i^*$.

\smallskip
\noindent
\textit{Part (b).} Let $\mu_i$ be stationary and average-reward optimal against $x_{-i}$. Define the nonnegative slack function
\[
d(s,a_i):=g_i^*+h_i(s)-r_i^{x_{-i}}(s,a_i)-\sum_{s'\in S}Q_i^{x_{-i}}(s'\mid s,a_i)h_i(s').
\]
By construction, $d(s,a_i)\ge0$ for every state-action pair. Let $\pi$ be the unique invariant distribution of the irreducible chain $P^{\mu_i,x_{-i}}$. Then
\begin{align*}
\sum_{s\in S}\pi(s)\sum_{a_i\in A_i(s)}\mu_i(a_i\mid s)d(s,a_i)
&=\sum_{s}\pi(s)\left(g_i^*+h_i(s)-\bar r_i^{\mu_i,x_{-i}}(s)-\sum_{s'}P^{\mu_i,x_{-i}}(s,s')h_i(s')\right)\\
&=g_i^*-\sum_s\pi(s)\bar r_i^{\mu_i,x_{-i}}(s)\\
&=g_i^*-g_i(\mu_i,x_{-i})\\
&=0.
\end{align*}
The second equality uses the invariance identity $\sum_s\pi(s)h_i(s)=\sum_{s,s'}\pi(s)P^{\mu_i,x_{-i}}(s,s')h_i(s')$, and the last equality uses optimality. Since all terms are nonnegative and, because the chain is irreducible, $\pi(s)>0$ for every state $s$, we must have
\[
\sum_{a_i}\mu_i(a_i\mid s)d(s,a_i)=0\qquad\text{for every }s\in S.
\]
Thus whenever $\mu_i(a_i\mid s)>0$, necessarily $d(s,a_i)=0$, which means $a_i\in A_i^*(s;x_{-i})$. This proves the support inclusion.
\end{proof}

\section{The average-reward meta-game}\label{sec:metagame}

\subsection{Definition and continuity of the payoff map}

For a stationary mixed profile $x\in X$, define the expected one-step payoff of player $i$ in state $s$ by
\[
\bar u_i^x(s):=\sum_{a\in A(s)}x(a\mid s)u_i(s,a).
\]
By Proposition~\ref{prop:mixed-irreducible}, the chain $P^x$ is irreducible. Let $\pi^x$ denote its unique invariant distribution. We define the average-reward meta-game payoff of player $i$ at $x$ by
\begin{equation}\label{eq:meta-payoff}
g_i(x):=\sum_{s\in S}\pi^x(s)\bar u_i^x(s).
\end{equation}
Because $P^x$ is irreducible, this quantity is equal to the unique average reward of player $i$ under the stationary profile $x$, independent of the initial state.

\begin{proposition}[Continuity of invariant distributions and average rewards]\label{prop:continuity}
Assume Condition~\ref{cond:A}. For every player $i$, the map
\[
x\longmapsto \pi^x\in \Delta(S)
\]
is continuous on $X$, and therefore the payoff map
\[
x\longmapsto g_i(x)
\]
is continuous on $X$.
\end{proposition}

\begin{proof}
The map $x\mapsto P^x$ is multilinear, hence continuous. Let $(x^m)_{m\ge1}$ be a sequence in $X$ converging to $x\in X$. Write $\pi^m:=\pi^{x^m}$. Since $\Delta(S)$ is compact, every subsequence of $(\pi^m)$ admits a further convergent subsequence. Let $\pi^{m_k}\to \bar\pi$.

For each $k$, the invariance identity is
\[
\pi^{m_k}P^{x^{m_k}}=\pi^{m_k}.
\]
Passing to the limit and using continuity of matrix multiplication gives
\[
\bar\pi P^x=\bar\pi.
\]
Thus $\bar\pi$ is an invariant distribution for $P^x$. By Proposition~\ref{prop:mixed-irreducible}, $P^x$ is irreducible, so its invariant distribution is unique. Hence $\bar\pi=\pi^x$.

We have shown that every convergent subsequence of $(\pi^m)$ converges to $\pi^x$. Since $\Delta(S)$ is compact, this implies the whole sequence converges: $\pi^m\to \pi^x$. Therefore $x\mapsto \pi^x$ is continuous.

The map $x\mapsto \bar u_i^x(s)$ is continuous for each state $s$ because it is multilinear in the stationary mixed strategies. Formula \eqref{eq:meta-payoff} then shows that $g_i$ is continuous as the composition of continuous operations.
\end{proof}

The next proposition records the ergodic interpretation of $g_i(x)$. The proof is standard but useful because it links the meta-game payoff to the actual long-run behavior of play.

\begin{proposition}[Ergodic theorem for stationary profiles]\label{prop:ergodic}
Assume Condition~\ref{cond:A}. Fix a stationary mixed profile $x\in X$, an initial state $s\in S$, and a player $i$. Then under $\Prob_s^x$,
\[
\frac1T\sum_{t=1}^{T}u_i(s_t,a_t)\longrightarrow g_i(x)
\qquad\text{almost surely and in }L^1.
\]
\end{proposition}

\begin{proof}
Because $P^x$ is a finite irreducible Markov chain, the ergodic theorem for finite irreducible chains implies that for every bounded function $f:S\to\R$,
\[
\frac1T\sum_{t=1}^{T}f(s_t)\longrightarrow\sum_{z\in S}\pi^x(z)f(z)
\qquad\text{almost surely and in }L^1.
\]
Apply this to $f=\bar u_i^x$. Then
\begin{equation}\label{eq:ergodic-expected-reward}
\frac1T\sum_{t=1}^{T}\bar u_i^x(s_t)\longrightarrow \sum_{z\in S}\pi^x(z)\bar u_i^x(z)=g_i(x)
\qquad\text{almost surely and in }L^1.
\end{equation}

Next define
\[
d_t:=u_i(s_t,a_t)-\bar u_i^x(s_t),\qquad t\ge1.
\]
Let
\[
\mathcal G_t:=\sigma(s_1,a_1,s_2,a_2,\dots,s_t,a_t,s_{t+1}),\qquad t\ge0,
\]
with the convention that $\mathcal G_0=\sigma(s_1)$. Then $d_t$ is $\mathcal G_t$-measurable and, because under the stationary profile $x$ the action at stage $t$ is drawn according to $x(\cdot\mid s_t)$,
\[
\E[d_t\mid \mathcal G_{t-1}]=0.
\]
Hence the partial sums $D_T:=\sum_{t=1}^{T}d_t$ form a martingale with bounded increments. By the strong law for bounded martingale differences,
\[
\frac{D_T}{T}\longrightarrow 0
\qquad\text{almost surely and in }L^1.
\]
Combining this with \eqref{eq:ergodic-expected-reward} yields
\[
\frac1T\sum_{t=1}^{T}u_i(s_t,a_t)
=\frac1T\sum_{t=1}^{T}\bar u_i^x(s_t)+\frac{D_T}{T}
\longrightarrow g_i(x)
\]
almost surely and in $L^1$.
\end{proof}

\subsection{Best responses in the meta-game}

Fix a player $i$. The best-response correspondence of player $i$ in the average-reward meta-game is
\[
B_i(x_{-i}):=\arg\max_{\mu_i\in X_i} g_i(\mu_i,x_{-i}),\qquad x_{-i}\in X_{-i}.
\]

\begin{proposition}[Structure of best responses]\label{prop:best-response-structure}
Assume Condition~\ref{cond:A}. Fix a player $i$ and $x_{-i}\in X_{-i}$. Then
\begin{equation}\label{eq:best-response-product}
B_i(x_{-i})=
\Bigl\{\mu_i\in X_i:\supp\mu_i(\cdot\mid s)\subseteq A_i^*(s;x_{-i})\text{ for every }s\in S\Bigr\}.
\end{equation}
In particular, $B_i(x_{-i})$ is nonempty, compact, and convex.
\end{proposition}

\begin{proof}
Because $X_i$ is compact and $\mu_i\mapsto g_i(\mu_i,x_{-i})$ is continuous by Proposition~\ref{prop:continuity}, the set $B_i(x_{-i})$ is nonempty and compact.

By Proposition~\ref{prop:optimal-stationary}, a stationary mixed strategy $\mu_i$ is average-reward optimal against $x_{-i}$ if and only if its support is contained statewise in the maximizing sets $A_i^*(s;x_{-i})$. This proves \eqref{eq:best-response-product}. The right-hand side is a product of simplices
\[
\prod_{s\in S}\Delta\bigl(A_i^*(s;x_{-i})\bigr),
\]
hence convex and compact.
\end{proof}

\begin{proposition}[Closed graph of the best-response correspondence]\label{prop:best-response-closed-graph}
Assume Condition~\ref{cond:A}. For each player $i$, the correspondence $B_i:X_{-i}\rightrightarrows X_i$ has closed graph. Consequently it is upper hemicontinuous.
\end{proposition}

\begin{proof}
Let $x_{-i}^m\to x_{-i}$ in $X_{-i}$, let $\mu_i^m\to \mu_i$ in $X_i$, and suppose $\mu_i^m\in B_i(x_{-i}^m)$ for every $m$. Then for every fixed $\eta_i\in X_i$,
\[
g_i(\mu_i^m,x_{-i}^m)\ge g_i(\eta_i,x_{-i}^m).
\]
Passing to the limit and using continuity of $g_i$ from Proposition~\ref{prop:continuity}, we get
\[
g_i(\mu_i,x_{-i})\ge g_i(\eta_i,x_{-i})\qquad\text{for every }\eta_i\in X_i.
\]
Hence $\mu_i\in B_i(x_{-i})$, so the graph is closed.

Because the domain $X_{-i}$ is compact and the values of $B_i$ are compact by Proposition~\ref{prop:best-response-structure}, closed graph implies upper hemicontinuity.
\end{proof}

\begin{remark}[Occupation-measure formulation]\label{rem:occupation}
For fixed $i$ and $x_{-i}$, the average-reward MDP can also be written as a linear program over state-action occupation measures. One maximizes
\[
\sum_{s\in S}\sum_{a_i\in A_i(s)}m(s,a_i)\,r_i^{x_{-i}}(s,a_i)
\]
subject to the balance equations
\[
\sum_{a_i}m(s,a_i)=\sum_{z\in S}\sum_{b_i\in A_i(z)}m(z,b_i)Q_i^{x_{-i}}(s\mid z,b_i),
\qquad s\in S,
\]
together with $m\ge0$ and $\sum_{s,a_i}m(s,a_i)=1$. The feasible set is convex, so the set of $\eta$-optimal \emph{occupation measures} is convex for every $\eta\ge0$. We will not use this formulation below, because Proposition~\ref{prop:best-response-structure} already gives the stronger policy-space description of exact best responses that is needed for the fixed-point argument.
\end{remark}

\subsection{Existence of a stationary average-reward equilibrium}

We now apply a standard fixed-point theorem to the product best-response correspondence.

\begin{theorem}[Kakutani-Fan-Glicksberg fixed-point theorem]\label{thm:kfg}
Let $K$ be a nonempty compact convex subset of a Euclidean space, and let $\Phi:K\rightrightarrows K$ be a correspondence with nonempty compact convex values and closed graph. Then there exists $x\in K$ such that $x\in \Phi(x)$.
\end{theorem}

\begin{proof}
This is the classical Kakutani-Fan-Glicksberg fixed-point theorem; see Glicksberg \cite{Glicksberg1952}.
\end{proof}

\begin{theorem}[Stationary equilibrium of the average-reward meta-game]\label{thm:meta-equilibrium}
Assume Condition~\ref{cond:A}. Then there exists a stationary mixed profile $x^*\in X$ such that for every player $i$,
\[
g_i(x^*)\ge g_i(\mu_i,x_{-i}^*)\qquad\text{for every }\mu_i\in X_i.
\]
Equivalently, $x_i^*\in B_i(x_{-i}^*)$ for every player $i$.
\end{theorem}

\begin{proof}
Define the product correspondence $B:X\rightrightarrows X$ by
\[
B(x):=\prod_{i\in N}B_i(x_{-i}).
\]
The set $X$ is a nonempty compact convex subset of a finite-dimensional Euclidean space. By Propositions~\ref{prop:best-response-structure} and \ref{prop:best-response-closed-graph}, each $B_i$ has nonempty compact convex values and closed graph. Therefore $B$ itself has nonempty compact convex values and closed graph.

By Theorem~\ref{thm:kfg}, there exists $x^*\in X$ with $x^*\in B(x^*)$. This means precisely that $x_i^*\in B_i(x_{-i}^*)$ for every player $i$, or equivalently that each $x_i^*$ maximizes $g_i(\mu_i,x_{-i}^*)$ over $\mu_i\in X_i$.
\end{proof}

\section{Proof of the main theorem}\label{sec:proof-main}

We can now prove the uniform-equilibrium theorem announced in the introduction.

\begin{proof}[Proof of Theorem~\ref{thm:intro-main}]
Fix $\varepsilon>0$. By Theorem~\ref{thm:meta-equilibrium}, choose a stationary mixed profile
\[
x^*=(x_1^*,\dots,x_n^*)\in X
\]
such that each $x_i^*$ is a best response to $x_{-i}^*$ in the average-reward meta-game.

Now fix a player $i$. Consider the average-reward MDP faced by player $i$ against the stationary mixed opponents' profile $x_{-i}^*$. Let
\[
g_i^*:=g_i^*(x_{-i}^*),\qquad h_i^*:=h_i^{x_{-i}^*},\qquad H_i:=\spn(h_i^*).
\]
Because $x_i^*\in B_i(x_{-i}^*)$, we have
\[
g_i(x^*)=\max_{\mu_i\in X_i}g_i(\mu_i,x_{-i}^*)=g_i^*.
\]
Moreover, by Proposition~\ref{prop:best-response-structure}, the support of $x_i^*(\cdot\mid s)$ is contained in the maximizing set $A_i^*(s;x_{-i}^*)$ for every state $s$.

Let $s\in S$ be any initial state, let $T\ge1$, and let $\tau_i$ be any behavioral deviation of player $i$.

\smallskip
\noindent
\textbf{Step 1: upper bound on the deviator.}
By Proposition~\ref{prop:supermartingale},
\begin{equation}\label{eq:proof-dev-upper}
\gamma_i^T\bigl(s,(\tau_i,x_{-i}^*)\bigr)\le g_i^*+\frac{H_i}{T}.
\end{equation}
This bound holds for \emph{every} behavioral deviation $\tau_i$, including arbitrary history-dependent randomized deviations.

\smallskip
\noindent
\textbf{Step 2: lower bound on the compliant payoff.}
Since $x_i^*$ mixes only over maximizing actions, Proposition~\ref{prop:optimal-stationary}(a) yields the exact identity
\[
\gamma_i^T(s,x^*)=g_i^*+\frac{h_i^*(s)-\E_s^{x^*}[h_i^*(s_{T+1})]}{T}.
\]
Hence
\begin{equation}\label{eq:proof-compliant-lower}
\gamma_i^T(s,x^*)\ge g_i^* - \frac{H_i}{T}.
\end{equation}
By Proposition~\ref{prop:ergodic}, one also has almost sure and $L^1$ convergence
\[
\frac1T\sum_{t=1}^{T}u_i(s_t,a_t)\longrightarrow g_i(x^*)=g_i^*
\qquad\text{under }\Prob_s^{x^*}.
\]
Thus the compliant path has the correct ergodic average, and \eqref{eq:proof-compliant-lower} gives the quantitative finite-horizon estimate needed for uniform equilibrium.

\smallskip
\noindent
\textbf{Step 3: regret bound.}
Subtracting \eqref{eq:proof-compliant-lower} from \eqref{eq:proof-dev-upper} gives
\[
\gamma_i^T\bigl(s,(\tau_i,x_{-i}^*)\bigr)-\gamma_i^T(s,x^*)\le \frac{2H_i}{T}.
\]
Since this holds for every player $i$, every initial state $s$, every horizon $T\ge1$, and every behavioral deviation $\tau_i$, we obtain the uniform estimate
\[
\sup_{i\in N}\sup_{s\in S}\sup_{\tau_i}\Bigl[\gamma_i^T\bigl(s,(\tau_i,x_{-i}^*)\bigr)-\gamma_i^T(s,x^*)\Bigr]
\le \frac{2H}{T},
\]
where
\[
H:=\max_{i\in N}H_i=\max_{i\in N}\spn(h_i^*).
\]

\smallskip
\noindent
\textbf{Step 4: choose the threshold.}
Set
\[
T_0:=\left\lceil \frac{2H}{\varepsilon}\right\rceil.
\]
Then for every $T\ge T_0$,
\[
\gamma_i^T(s,x^*)\ge \gamma_i^T\bigl(s,(\tau_i,x_{-i}^*)\bigr)-\varepsilon
\]
for every player $i$, every initial state $s$, and every behavioral deviation $\tau_i$.

Therefore $x^*$ is a uniform $\varepsilon$-equilibrium. This proves Theorem~\ref{thm:intro-main}.
\end{proof}

For later reference, we record the quantitative consequence explicitly.

\begin{corollary}[Uniform $O(1/T)$ regret]\label{cor:o1overT}
Under Condition~\ref{cond:A}, the stationary profile $x^*$ from Theorem~\ref{thm:meta-equilibrium} satisfies
\[
\gamma_i^T\bigl(s,(\tau_i,x_{-i}^*)\bigr)-\gamma_i^T(s,x^*)
\le \frac{2\spn(h_i^*)}{T}
\le \frac{2\max_j\spn(h_j^*)}{T}
\]
for every player $i$, every initial state $s$, every horizon $T\ge1$, and every behavioral deviation $\tau_i$.
\end{corollary}

\section{Discussion}\label{sec:discussion}

The proof above is intentionally global. Condition~\ref{cond:A} is a hypothesis about the full Markov chain on the original state space $S$ under unilateral pure stationary deviations. Once that condition is imposed at the game level, every stationary mixed profile is irreducible, every unilateral stationary control problem is a unichain average-reward MDP on $S$, and the whole argument can be written directly on the original game.

Three points are worth emphasizing.

\medskip
\noindent
\textbf{1. No decomposition is needed.}
Under Condition~\ref{cond:A}, the right state-space decomposition is the trivial one: the whole game is one communicating object. Any auxiliary expansion of the state space, such as a routing tree or an added controller memory, may create recurrent classes that do not exist in the original game. The present proof avoids that entirely.

\medskip
\noindent
\textbf{2. The equilibrium step is a compact fixed-point argument on stationary mixed strategies.}
The meta-game payoff map $g$ is continuous because the invariant distribution of a finite irreducible chain depends continuously on the transition matrix. The fixed-point argument does not require any local Bellman certificates, tree packages, or auxiliary balancing equations. What matters is simply that best responses are average-reward optimal stationary policies of unichain MDPs, and these best-response sets are products of simplices over maximizing actions.

\medskip
\noindent
\textbf{3. The deviation analysis is genuinely policy-uniform.}
The supermartingale in Proposition~\ref{prop:supermartingale} does not care whether the deviator is stationary, Markov, finite-memory, or fully history-dependent. The only ingredient is the Bellman inequality for the bias function, which is pointwise in the current state and action. This is why the deviation bound is so clean: once the opponents are stationary, every unilateral deviation sees an MDP, and the bias function controls all of them at once.

The theorem leaves the general multiplayer problem untouched. Without Condition~\ref{cond:A}, a unilateral decision problem can be multichain, the gain can depend on the recurrent class selected by the policy, and a single global bias no longer controls finite-horizon deviations uniformly across initial states. That multichain obstacle is precisely what the present hypothesis rules out. The general existence problem for uniform equilibria in finite stochastic games remains open \cite{SolanVieille2015,Solan2022,OpenProblems2026}.

Still, the result is useful in its own right. It yields a stationary uniform $\varepsilon$-equilibrium with explicit $O(1/T)$ regret under a transparent and checkable game-level irreducibility condition. Just as importantly, it shows that under this condition the proof really is a one-node argument: average-reward MDP theory plus a single fixed-point theorem are enough.

\begin{thebibliography}{99}

\bibitem{Glicksberg1952}
I.~L. Glicksberg.
\newblock A further generalization of the Kakutani fixed point theorem, with application to Nash equilibrium points.
\newblock \emph{Proceedings of the American Mathematical Society}, 3(1):170--174, 1952.

\bibitem{MertensNeyman1981}
J.-F. Mertens and A. Neyman.
\newblock Stochastic games.
\newblock \emph{International Journal of Game Theory}, 10(2):53--66, 1981.

\bibitem{Neyman2003Minmax}
A. Neyman.
\newblock Stochastic games: existence of the minmax.
\newblock In A. Neyman and S. Sorin, editors, \emph{Stochastic Games and Applications}, volume 570 of \emph{NATO Science Series C: Mathematical and Physical Sciences}, pages 173--193. Springer, Dordrecht, 2003.

\bibitem{Neyman2003MarkovChains}
A. Neyman.
\newblock From Markov chains to stochastic games.
\newblock In A. Neyman and S. Sorin, editors, \emph{Stochastic Games and Applications}, volume 570 of \emph{NATO Science Series C: Mathematical and Physical Sciences}, pages 9--25. Springer, Dordrecht, 2003.

\bibitem{Norris1998}
J.~R. Norris.
\newblock \emph{Markov Chains}.
\newblock Cambridge University Press, Cambridge, 1998.

\bibitem{OpenProblems2026}
Game Theory Society (Israel Chapter).
\newblock Open problems.
\newblock \url{https://sites.google.com/view/thegametheorysocietyilchapter/open-problems}, accessed April 2, 2026.

\bibitem{Puterman1994}
M.~L. Puterman.
\newblock \emph{Markov Decision Processes: Discrete Stochastic Dynamic Programming}.
\newblock Wiley, New York, 1994.

\bibitem{Shapley1953}
L.~S. Shapley.
\newblock Stochastic games.
\newblock \emph{Proceedings of the National Academy of Sciences of the United States of America}, 39(10):1095--1100, 1953.

\bibitem{Solan2022}
E. Solan.
\newblock \emph{A Course in Stochastic Game Theory}.
\newblock Cambridge University Press, Cambridge, 2022.

\bibitem{SolanVieille2015}
E. Solan and N. Vieille.
\newblock Stochastic games.
\newblock \emph{Proceedings of the National Academy of Sciences of the United States of America}, 112(45):13743--13746, 2015.

\bibitem{Vieille2000a}
N. Vieille.
\newblock Two-player stochastic games {I}: A reduction.
\newblock \emph{Israel Journal of Mathematics}, 119:55--91, 2000.

\bibitem{Vieille2000b}
N. Vieille.
\newblock Two-player stochastic games {II}: The case of recursive games.
\newblock \emph{Israel Journal of Mathematics}, 119:93--126, 2000.

\end{thebibliography}

\end{document}
